\section{Tensor products}
\label{section-tensor-product}

\begin{definition}
\label{definition-bilinear}
Let $R$ be a ring, $M, N, P$ be three $R$-modules.
A mapping $f : M \times N \to P$ (where $M \times N$
is viewed only as Cartesian product of two $R$-modules) is said to be
{\it $R$-bilinear} if for each $x \in M$
the mapping $y\mapsto f(x, y)$ of $N$ into $P$ is $R$-linear, and for each
$y\in N$ the mapping $x\mapsto f(x, y)$ is also $R$-linear.
\end{definition}

\begin{lemma}
\label{lemma-tensor-product}
Let $M, N$ be $R$-modules. Then there exists a pair $(T, g)$
where $T$ is an $R$-module, and
$g : M \times N \to T$ an $R$-bilinear
mapping, with the following universal property:
For any $R$-module $P$ and any $R$-bilinear mapping
$f : M \times N \to P$, there
exists a unique $R$-linear
mapping $\tilde{f} : T \to P$ such that $f = \tilde{f} \circ g$.
In other words, the following diagram commutes:
$$
\xymatrix{
M \times N \ar[rr]^f \ar[dr]_g & & P\\
& T \ar[ur]_{\tilde f}
}
$$
Moreover, if $(T, g)$ and $(T', g')$
are two pairs with this property, then there
exists a unique isomorphism
$j : T \to T'$ such that $j\circ g = g'$.
\end{lemma}

\noindent
The $R$-module $T$ which satisfies the above universal property is called
the  {\it tensor product} of $R$-modules $M$ and $N$, denoted as
$M \otimes_R N$.

\begin{proof}
We first prove the existence of such $R$-module $T$.
Let $M, N$ be $R$-modules.
Let $T$ be the quotient module
$P/Q$, where $P$ is the free $R$-module $R^{(M \times N)}$ and $Q$ is the
$R$-module generated by all elements of
the following types: ($x\in M, y\in N$)
\begin{align*}
(x + x', y) - (x, y) - (x', y), \\
(x, y + y') - (x, y) - (x, y'), \\
(ax, y) - a(x, y), \\
(x, ay) - a(x, y)
\end{align*}
Let $\pi : M \times N \to T$ denote the natural map.
This map is $R$-bilinear, as
implied by the above relations
when we check the bilinearity conditions. Denote the image
$\pi(x, y) = x \otimes
y$, then these elements generate
$T$. Now let $f : M \times N \to P$ be an $R$-bilinear map,
then we can define
$f' : T \to P$ by extending the mapping
$f'(x \otimes y) = f(x, y)$. Clearly $f = f'\circ \pi$. Moreover, $f'$ is
uniquely determined by the value on the
generating sets $\{x \otimes y : x\in M, y\in N\}$.
Suppose there is another pair $(T', g')$ satisfying the same properties.
Then there is a unique $j : T \to T'$ and
also $j' : T' \to T$ such that $g' = j\circ g$, $g = j'\circ g'$.
But then both the maps $(j\circ j') \circ g$ and $g$
satisfies the universal properties, so by uniqueness they are equal,
and hence $j'\circ j$ is identity on $T$.
Similarly $(j'\circ j) \circ g' = g'$ and $j\circ j'$ is identity on $T'$.
So $j$ is an isomorphism.
\end{proof}

\begin{lemma}
\label{lemma-flip-tensor-product}
Let $M, N, P$ be $R$-modules, then the bilinear maps
\begin{align*}
(x, y) & \mapsto y \otimes x\\
(x + y, z) & \mapsto x \otimes z + y \otimes z\\
(r, x) & \mapsto rx
\end{align*}
induce unique isomorphisms
\begin{align*}
M \otimes_R N & \to N \otimes_R M, \\
(M\oplus N)\otimes_R P & \to (M \otimes_R P)\oplus(N \otimes_R P),  \\
R \otimes_R M & \to M
\end{align*}
\end{lemma}

\begin{proof}
Omitted.
\end{proof}

\noindent
We may generalize the tensor product of two $R$-modules to finitely many
$R$-modules, and set up a
correspondence between the multi-tensor product with multilinear mappings.
Using almost the same construction
one can prove that:

\begin{lemma}
\label{lemma-multilinear}
Let $M_1, \ldots, M_r$ be $R$-modules. Then there exists a pair $(T, g)$
consisting of an $R$-module T and an $R$-multilinear mapping
$g : M_1\times \ldots \times M_r \to T$ with the universal
property: For any $R$-multilinear mapping
$f : M_1\times \ldots \times M_r \to P$ there exists a unique $R$-module
homomorphism $f' : T \to P$ such that $f'\circ g = f$.
Such a module $T$ is unique up to unique isomorphism. We denote it
$M_1\otimes_R \ldots \otimes_R M_r$ and we denote the universal
multilinear map $(m_1, \ldots, m_r) \mapsto m_1 \otimes \ldots \otimes m_r$.
\end{lemma}

\begin{proof}
Omitted.
\end{proof}

\begin{lemma}
\label{lemma-transitive}
The homomorphisms
$$
(M \otimes_R N)\otimes_R P \to
M \otimes_R N \otimes_R P \to
M \otimes_R (N \otimes_R P)
$$
such that
$f((x \otimes y)\otimes z) = x \otimes y \otimes z$
and $g(x \otimes y \otimes z) = x \otimes (y \otimes z)$,
$x\in M, y\in N, z\in P$ are well-defined and are isomorphisms.
\end{lemma}

\begin{proof}
We shall prove $f$ is well-defined and is an isomorphism, and this proof
carries analogously to $g$. Fix any
$z\in P$, then the mapping $(x, y)\mapsto x \otimes y \otimes z$,
$x\in M, y\in N$, is $R$-bilinear in $x$ and $y$,
and hence induces homomorphism $f_z : M \otimes N \to M \otimes N \otimes P$
which sends
$f_z(x \otimes y) = x \otimes y \otimes z$.
Then consider $(M \otimes N)\times P \to M \otimes N \otimes P$ given by
$(w, z)\mapsto f_z(w)$. The map is
$R$-bilinear and thus induces
$f : (M \otimes_R N)\otimes_R P \to M \otimes_R N \otimes_R P$
and $f((x \otimes y)\otimes z) = x \otimes y \otimes z$.
To construct the inverse, we note that the map
$\pi : M \times N \times P \to (M \otimes N)\otimes P$ is
$R$-trilinear.
Therefore, it induces an $R$-linear map
$h : M \otimes N \otimes P \to (M \otimes N)\otimes P$ which
agrees with the universal property. Here we see that
$h(x \otimes y \otimes z) = (x \otimes y)\otimes z$.
From the explicit expression of $f$ and $h$, $f\circ h$ and $h\circ f$ are
identity maps of $M \otimes N \otimes
P$ and $(M \otimes N)\otimes P$ respectively, hence $f$ is our desired
isomorphism.
\end{proof}

\noindent
Doing induction we see that this extends to multi-tensor products. Combined
with Lemma \ref{lemma-flip-tensor-product} we see that
the tensor product operation on the category of $R$-modules is associative,
commutative and distributive.

\begin{definition}
\label{definition-bimodule}
An abelian group $N$ is called an {\it $(A, B)$-bimodule} if it is both an
$A$-module and a $B$-module and for all $a \in A$ and $b \in B$ the
multiplication by $a$ and $b$ commute, so $b(an) = a(bn)$ for all $n \in N$.
In this situation we usually write the $B$-action on the right: so
for $b \in B$ and $n \in N$ the result of multiplying $n$ by $b$
is denoted $nb$. With this convention the compatibility above is
that $(ax)b = a(xb)$ for all $a\in A, b\in B, x\in N$.
The shorthand $_AN_B$ is used to denote an $(A, B)$-bimodule $N$.
\end{definition}

\begin{lemma}
\label{lemma-tensor-with-bimodule}
For $A$-module $M$, $B$-module $P$ and $(A, B)$-bimodule $N$, the modules
$(M \otimes_A N)\otimes_B P$ and $M \otimes_A(N \otimes_B P)$ can both be
given $(A, B)$-bimodule structure,
and moreover
$$
(M \otimes_A N)\otimes_B P \cong M \otimes_A(N \otimes_B P).
$$
\end{lemma}

\begin{proof}
A priori $M \otimes_A N$ is an $A$-module, but we can give it a
$B$-module structure by letting
$$
(x \otimes y)b = x \otimes yb, \quad x\in M, y\in N, b\in B
$$
Thus $M \otimes_A N$ becomes an $(A, B)$-bimodule. Similarly for
$N \otimes_B P$, and thus for
$(M \otimes_A N)\otimes_B P$ and $M \otimes_A(N \otimes_B P)$. By
Lemma \ref{lemma-transitive}, these two
modules are isomorphic as both as $A$-module and $B$-module via the same
mapping.
\end{proof}

\begin{lemma}
\label{lemma-hom-from-tensor-product}
For any three $R$-modules $M, N, P$,
$$
\Hom_R(M \otimes_R N, P) \cong \Hom_R(M, \Hom_R(N, P))
$$
\end{lemma}

\begin{proof}
An $R$-linear map $\hat{f}\in \Hom_R(M \otimes_R N, P)$ corresponds to an
$R$-bilinear map $f : M \times N \to P$. For
each $x\in M$ the mapping $y\mapsto f(x, y)$ is $R$-linear by the universal
property. Thus $f$ corresponds to a
map $\phi_f : M \to \Hom_R(N, P)$. This map is $R$-linear since
$$
\phi_f(ax + y)(z) =
f(ax + y, z) = af(x, z)+f(y, z) =
(a\phi_f(x)+\phi_f(y))(z),
$$
for all $a \in R$, $x \in M$, $y \in M$ and
$z \in N$. Conversely, any
$f \in \Hom_R(M, \Hom_R(N, P))$ defines an $R$-bilinear
map $M \times N \to P$, namely $(x, y)\mapsto f(x)(y)$.
So this is a natural one-to-one correspondence between the
two modules
$\Hom_R(M \otimes_R N, P)$ and $\Hom_R(M, \Hom_R(N, P))$.
\end{proof}

\begin{lemma}[Tensor products commute with colimits]
\label{lemma-tensor-products-commute-with-limits}
Let $(M_i, \mu_{ij})$ be a system over the preordered set $I$.
Let $N$ be an $R$-module. Then
$$
\colim (M_i \otimes N) \cong (\colim M_i)\otimes N.
$$
Moreover, the isomorphism is induced by the homomorphisms
$\mu_i \otimes 1: M_i \otimes N \to M \otimes N$
where $M = \colim_i M_i$ with natural maps $\mu_i : M_i \to M$.
\end{lemma}

\begin{proof}
First proof. The functor $M' \mapsto M' \otimes_R N$ is left adjoint
to the functor $N' \mapsto \Hom_R(N, N')$ by
Lemma \ref{lemma-hom-from-tensor-product}. Thus $M' \mapsto M' \otimes_R N$
commutes with all colimits, see
Categories, Lemma \ref{categories-lemma-adjoint-exact}.

\medskip\noindent
Second direct proof. Let $P = \colim (M_i \otimes N)$ with coprojections
$\lambda_i : M_i \otimes N \to P$. Let $M = \colim M_i$ with coprojections
$\mu_i : M_i \to M$. Then for all $i\leq j$, the following diagram commutes:
$$
\xymatrix{
M_i \otimes N \ar[r]_{\mu_i \otimes 1} \ar[d]_{\mu_{ij} \otimes 1} &
M \otimes N \ar[d]^{\text{id}} \\
M_j \otimes N \ar[r]^{\mu_j \otimes 1} &
M \otimes N
}
$$
By Lemma \ref{lemma-homomorphism-limit} these maps induce a unique homomorphism
$\psi : P \to M \otimes N$ such that
$\mu_i \otimes 1 = \psi \circ \lambda_i$.

\medskip\noindent
To construct the inverse map, for each $i\in I$, there is the canonical
$R$-bilinear mapping $g_i : M_i \times N \to
M_i \otimes N$. This induces a unique mapping
$\widehat{\phi} : M \times N \to P$
such that $\widehat{\phi} \circ (\mu_i \times 1) = \lambda_i \circ g_i$.
It is $R$-bilinear. Thus it induces an
$R$-linear mapping $\phi : M \otimes N \to P$.
From the commutative diagram below:
$$
\xymatrix{
M_i \times N \ar[r]^{g_i} \ar[d]^{\mu_i \times \text{id}} &
M_i \otimes N\ar[r]_{\text{id}} \ar[d]_{\lambda_i} &
M_i \otimes N \ar[d]_{\mu_i \otimes \text{id}} \ar[rd]^{\lambda_i} \\
M \times N \ar[r]^{\widehat{\phi}} &
P \ar[r]^{\psi} & M \otimes N \ar[r]^{\phi} & P
}
$$
we see that $\psi\circ\widehat{\phi} = g$, the canonical $R$-bilinear mapping
$g : M \times N \to M \otimes N$. So
$\psi\circ\phi$ is identity on $M \otimes N$. From the right-hand square and
triangle, $\phi\circ\psi$ is also
identity on $P$.
\end{proof}

\begin{lemma}
\label{lemma-tensor-product-exact}
Let
\begin{align*}
M_1\xrightarrow{f} M_2\xrightarrow{g} M_3 \to 0
\end{align*}
be an exact sequence of $R$-modules and homomorphisms, and let $N$ be any
$R$-module. Then the sequence
\begin{equation}
\label{equation-2ndex}
M_1\otimes N\xrightarrow{f \otimes 1} M_2\otimes N \xrightarrow{g \otimes 1}
M_3\otimes N \to 0
\end{equation}
is exact. In other words, the functor $- \otimes_R N$ is
{\it right exact}, in the sense that tensoring
each term in the original right exact sequence preserves the exactness.
\end{lemma}

\begin{proof}
For every $R$-module $P$
we apply the functor $\Hom(-, \Hom(N, P))$ to the first exact
sequence. We obtain
$$
0 \to
\Hom(M_3, \Hom(N, P)) \to
\Hom(M_2, \Hom(N, P)) \to
\Hom(M_1, \Hom(N, P))
$$
which is exact by Lemma \ref{lemma-hom-exact} (1).
By Lemma \ref{lemma-hom-from-tensor-product} this becomes the sequence
$$
0 \to \Hom(M_3 \otimes N, P) \to
\Hom(M_2 \otimes N, P) \to \Hom(M_1 \otimes N, P)
$$
which is therefore also exact. Then using
Lemma \ref{lemma-hom-exact} (1) again, we arrive at the desired exact sequence.
\end{proof}

\begin{remark}
\label{remark-tensor-product-not-exact}
However, tensor product does NOT preserve exact sequences in general.
In other words, if $M_1 \to M_2 \to M_3$ is
exact, then it is not necessarily true that
$M_1 \otimes N \to M_2 \otimes N \to M_3 \otimes N$
is exact for arbitrary $R$-module $N$.
\end{remark}

\begin{example}
\label{example-tensor-product-not-exact}
Consider the injective map $2 : \mathbf{Z}\to \mathbf{Z}$
viewed as a map of $\mathbf{Z}$-modules.
Let $N = \mathbf{Z}/2$. Then the induced map
$\mathbf{Z} \otimes \mathbf{Z}/2 \to \mathbf{Z} \otimes \mathbf{Z}/2$
is NOT injective. This is because for
$x \otimes y\in \mathbf{Z} \otimes \mathbf{Z}/2$,
$$
(2 \otimes 1)(x \otimes y) = 2x \otimes y = x \otimes 2y = x \otimes 0 = 0
$$
Therefore the induced map is the zero map while $\mathbf{Z} \otimes N\neq 0$.
\end{example}

\begin{remark}
\label{remark-flat-module}
For $R$-modules $N$, if the
functor $-\otimes_R N$ is exact, i.e. tensoring
with $N$ preserves all exact
sequences, then $N$ is said to be {\it flat} $R$-module.
We will discuss this later in Section \ref{section-flat}.
\end{remark}

\begin{lemma}
\label{lemma-tensor-finiteness}
Let $R$ be a ring. Let $M$ and $N$ be $R$-modules.
\begin{enumerate}
\item If $N$ and $M$ are finite, then so is $M \otimes_R N$.
\item If $N$ and $M$ are finitely presented, then so is $M \otimes_R N$.
\end{enumerate}
\end{lemma}

\begin{proof}
Suppose $M$ is finite. Then choose a presentation
$0 \to K \to R^{\oplus n} \to M \to 0$. This gives an exact sequence
$K \otimes_R N \to N^{\oplus n} \to M \otimes_R N \to 0$ by
Lemma \ref{lemma-tensor-product-exact}.
We conclude that if $N$ is finite too then $M \otimes_R N$
is a quotient of a finite module, hence finite, see
Lemma \ref{lemma-extension}.
Similarly, if both $N$ and $M$ are finitely presented, then
we see that $K$ is finite and that $M \otimes_R N$
is a quotient of the finitely presented module $N^{\oplus n}$ by
a finite module, namely $K \otimes_R N$, and hence finitely presented, see
Lemma \ref{lemma-extension}.
\end{proof}

\begin{lemma}
\label{lemma-tensor-localization}
Let $M$ be an $R$-module. Then the $S^{-1}R$-modules $S^{-1}M$
and $S^{-1}R \otimes_R M$ are canonically isomorphic, and the
canonical isomorphism $f : S^{-1}R \otimes_R M \to S^{-1}M$
is given by
$$
f((a/s) \otimes m) = am/s, \forall a \in R, m \in M, s \in S
$$
\end{lemma}

\begin{proof}
Obviously, the map
$f' : S^{-1}R \times M \to S^{-1}M$ given by $f'(a/s, m) = am/s$ is
bilinear, and thus by the
universal property, this map induces a unique $S^{-1}R$-module homomorphism
$f : S^{-1}R \otimes_R M \to S^{-1}M$ as in the statement of the lemma.
Actually every element in $S^{-1}M$ is of the form $m/s$, $m\in M, s\in S$ and
every element in
$S^{-1}R \otimes_R M$ is of the form $1/s \otimes m$. To see the latter fact,
write an element in
$S^{-1}R \otimes_R M$ as
$$
\sum_k \frac{a_k}{s_k} \otimes m_k =
\sum_k \frac{a_k t_k}{s} \otimes m_k =
\frac{1}{s} \otimes \sum_k {a_k t_k}m_k = \frac{1}{s} \otimes m
$$
Where $m = \sum_k {a_k t_k}m_k$. Then it is obvious that $f$ is surjective,
and if $f(\frac{1}{s} \otimes m) = m/s = 0$ then there exists $t \in S$ with
$tm = 0$ in $M$. Then we have
$$
\frac{1}{s} \otimes m = \frac{1}{st} \otimes tm = \frac{1}{st} \otimes 0 = 0
$$
Therefore $f$ is injective.
\end{proof}

\begin{lemma}
\label{lemma-tensor-product-localization}
Let $M, N$ be $R$-modules, then there is a canonical
$S^{-1}R$-module isomorphism
$f : S^{-1}M \otimes_{S^{-1}R}S^{-1}N \to S^{-1}(M \otimes_R N)$,
given by
$$
f((m/s)\otimes(n/t)) = (m \otimes n)/st
$$
\end{lemma}

\begin{proof}
We may use Lemma \ref{lemma-tensor-with-bimodule}
and Lemma \ref{lemma-tensor-localization} repeatedly to
see that these two
$S^{-1}R$-modules are isomorphic, noting that $S^{-1}R$ is an
$(R, S^{-1}R)$-bimodule:
\begin{align*}
S^{-1}(M \otimes_R N) & \cong S^{-1}R \otimes_R (M \otimes_R N)\\
 & \cong S^{-1}M \otimes_R N\\
 & \cong (S^{-1}M \otimes_{S^{-1}R}S^{-1}R)\otimes_R N\\
 & \cong S^{-1}M \otimes_{S^{-1}R}(S^{-1}R \otimes_R N)\\
 & \cong S^{-1}M \otimes_{S^{-1}R}S^{-1}N
\end{align*}
This isomorphism is easily seen to be the one stated in the lemma.
\end{proof}















